\section*{अध्याय \ref{ch:b2:mass-spec-atomic} --- द्रव्यमान स्पेक्ट्रममिति और परमाण्वीय स्पेक्ट्रमिकी}

\begin{solution}{exo:b2:mass-spec-atomic:1}
78; $400/2 = 200$; $194 + 1 = 195$, $[\mathrm M + \mathrm H]^+$।
\end{solution}

\begin{solution}{exo:b2:mass-spec-atomic:2}
विषम द्रव्यमान: नाइट्रोजन परमाणुओं की विषम संख्या, इतने छोटे अणु के लिए एक। \ce{C3H7NO} (प्रोपेनैमाइड, 73)
या \ce{C4H11N} (ब्यूटेन-1-ऐमीन, 73)।
\end{solution}

\begin{solution}{exo:b2:mass-spec-atomic:3}
लगभग समान: एक ब्रोमीन। 3 : 1: एक क्लोरीन।
\end{solution}

\begin{solution}{exo:b2:mass-spec-atomic:4}
पीला: सोडियम, \qty{589.0}{nm} और \qty{589.6}{nm}। किरमिजी: लिथियम, \qty{670.8}{nm}।
% ledger: asd:NaD2, asd:NaD1, asd:Li671
\end{solution}

\begin{solution}{exo:b2:mass-spec-atomic:5}
$n_C \approx 6.7/1.11 \approx 6$।
\end{solution}

\begin{solution}{exo:b2:mass-spec-atomic:6}
\ce{CO}: 27.9949; \ce{N2}: 28.0061; \ce{C2H4}: 28.0313 u। विभेदन क्षमताएँ: \ce{CO}/\ce{N2} $28/0.0112
\approx 2.5 \times 10^3$; \ce{N2}/\ce{C2H4} $28/0.0252 \approx 1.1 \times 10^3$; \ce{CO}/\ce{C2H4}
$28/0.0364 \approx 7.7 \times 10^2$।
% ledger: mass:C12, mass:O16, mass:N14, mass:H1
\end{solution}

\begin{solution}{exo:b2:mass-spec-atomic:7}
86: \ce{C5H10O} का \omterm{def:b2:mass-spec-atomic:molecular-ion}{आण्विक आयन}। 57: \ce{C2H5CO+}, $\alpha$-विदलन से बना एक \omterm{def:b2:aromatic-substitution:friedel-crafts}{ऐसिलियम आयन} (एक एथिल
मूलक की हानि, 29)। 29: \ce{C2H5+}।
% ledger: ms:pentan-3-one
\end{solution}

\begin{solution}{exo:b2:mass-spec-atomic:8}
86: $\mathrm M^{+\bullet}$। 71: $\mathrm M - \ce{CH3}$, ऐसिलियम \ce{C3H7CO+}। 43: \ce{CH3CO+}
(प्रोपिल मूलक की हानि)। 58: \omterm{def:b2:mass-spec-atomic:fragmentation}{मैकलैफ़र्टी पुनर्विन्यास}, प्रोपिल शृंखला के
$\gamma$ कार्बन के एक हाइड्रोजन से होकर, एथीन की हानि। पेंटेन-3-ओन में एथिल समूहों का कोई $\gamma$ कार्बन नहीं है: कोई
मैकलैफ़र्टी आयन नहीं; 58 पर उसका छोटा शिखर 57 वाले आयन का \ce{^{13}C} समस्थानिक शिखर है (तीन कार्बन,
लगभग 3.3~\%)।
% ledger: ms:pentan-2-one, ms:pentan-3-one
\end{solution}

\begin{solution}{exo:b2:mass-spec-atomic:9}
मूल-बिंदु से होकर एक रेखा: ढाल $\sum c_iA_i/\sum c_i^2 = 1.553/30.0 \approx \qty{0.0518}{L/mg}$। नमूना:
$0.128/0.0518 \approx \qty{2.47}{mg/L}$ कॉपर।
\end{solution}

\begin{solution}{exo:b2:mass-spec-atomic:10}\emergencystretch=3em
$t = L\sqrt{m/(2eU)} = 1.00\sqrt{1000 \times 1.6605 \times 10^{-27}/(2 \times 1.6022 \times 10^{-19}
\times 2.00 \times 10^4)} \approx \qty{16.1}{\micro s}$। 1001 u के लिए $t$ गुणक
$\sqrt{1.001} \approx 1 + 0.0005$ से बढ़ता है: $\Delta t \approx \qty{8.0}{ns}$।
% ledger: const:u, const:e
\end{solution}

\begin{solution}{exo:b2:mass-spec-atomic:11}
दो क्लोरीन: $(0.758 + 0.242x)^2$ से $0.575$, $0.367$, $0.0586$ मिलते हैं: $m/z$ 84, 86 और 88, अनुपात
100 : 64 : 10।
\end{solution}

\begin{solution}{exo:b2:mass-spec-atomic:12}
\ce{^{40}Ar^{16}O+}: $39.9624 + 15.9949 = 55.9573$; \ce{^{56}Fe}: 55.9349; $R = 56/0.0224 \approx 2.5 \times
10^3$। अन्यथा आयरन का कोई अन्य समस्थानिक, \ce{^{57}Fe}, मापिए, या विश्लेषक से पहले किसी संघट्ट सेल में
\ce{ArO+} को हटाइए।
% ledger: mass:Ar40, mass:Fe56, mass:O16
\end{solution}

\begin{solution}{pb:b2:mass-spec-atomic:1}
\textbf{1.} 156, 158 और 160 (157 और 159 के साथ): अणु कई समस्थानिक रूपों में होता है।
\textbf{2.} दो इकाई की दूरी पर तीन शिखर, बीच वाला सबसे बड़ा: एक क्लोरीन और एक ब्रोमीन।
\textbf{3.} 156।
\textbf{4.} सम द्रव्यमान: कोई नाइट्रोजन नहीं, या सम संख्या।
\textbf{5.} $156 - 79 - 35 = 42$: \ce{C3H6}।
\textbf{6.} \ce{C3H6BrCl}; असंतृप्ति की मात्रा $(2 \times 3 + 2 - 6 - 2)/2 = 0$: न कोई वलय, न कोई द्वि-आबंध।
\textbf{7.} $0.5/14.5 \approx 3.4~\%$, $3.4/1.11 \approx 3$ कार्बन।
\textbf{8.} 157 वाला शिखर 0.1 तक दिया गया है ऐसे पैमाने पर जहाँ वह 0.5 है: 20~\% अनिश्चितता।
\textbf{9.} \ce{C3H6^{81}Br^{35}Cl} और \ce{C3H6^{79}Br^{37}Cl}।
\textbf{10.} $3 \times 12 + 6 \times 1.007825 + 78.918338 + 34.968853 \approx 155.9341$।
\textbf{11.} $156/(156.151 - 155.934) \approx 7.2 \times 10^2$।
\textbf{12.} हाँ: दो हाइड्रोजन कम, 154।
\textbf{13.} एक ब्रोमीन परमाणु (79 या 81) खोया जाता है: \ce{C3H6Cl+}, \ce{^{35}Cl} के साथ 77, \ce{^{37}Cl} के साथ 79,
3 : 1।
\textbf{14.} HBr की हानि: मूलक धनायन $\mathrm{C_3H_5Cl^{+\bullet}}$, सम द्रव्यमान।
\textbf{15.} \ce{C3H5+}, ऐलिल धनायन, दोनों हैलोजनों की हानि के बाद (Br, फिर HCl)।
\textbf{16.} \ce{CH2Cl+} (49 और 51, 3 : 1)।
\textbf{17.} \ce{CH2Br+} (93 और 95) और \ce{C2H4Br+} (107 और 109), हर एक 1 : 1: शृंखला के एक सिरे पर ब्रोमीन,
दूसरा \ce{CH2Cl} की हानि से।
\textbf{18.} \ce{C-Br} आबंध \ce{C-Cl} आबंध से दुर्बल है, और \ce{Br} बेहतर निर्गामी
मूलक है।
\textbf{19.} नहीं: 127--131 पर कोई शिखर नहीं; \ce{CH2Cl+} और \ce{CH2Br+} दिखाते हैं कि हर हैलोजन एक अंतस्थ
\ce{CH2} पर है।
\textbf{20.} \ce{BrCH2CH2CH2Cl}, 1-ब्रोमो-3-क्लोरोप्रोपेन।
\textbf{21.} 156/158/160: M; 77/79: M $-$ Br; 76: M $-$ HBr; 107/109: M $-$ \ce{CH2Cl};
93/95: \ce{CH2Br+}; 49/51: \ce{CH2Cl+}; 41: \ce{C3H5+}; 39, 27: \ce{C3H3+}, \ce{C2H3+}।
\textbf{22.} नहीं: कोई कार्बोनिल समूह नहीं है।
\textbf{23.} 1-ब्रोमो-2-क्लोरोप्रोपेन, \ce{CH3CHClCH2Br}: \ce{CH2Br} की हानि 63/65 पर \ce{CH3CHCl+} देती है।
\textbf{24.} \ce{C3H6^{81}Br^{37}Cl}।
\textbf{25.} $M$: $0.758 \times 0.5069 = 0.384$; $M+2$: $0.758 \times 0.4931 + 0.242 \times 0.5069 =
0.496$; $M+4$: $0.242 \times 0.4931 = 0.119$। अनुपात $\boldsymbol{\approx 3.2 : 4.2 : 1}$, लगभग 3 : 4 :
1; स्पेक्ट्रम $14.5 : 18.6 : 4.4 = 3.3 : 4.2 : 1$ देता है।
\end{solution}
